% ============================================================
% The Character Structure of the Collision Fluctuation
% ============================================================

\documentclass[12pt]{amsart}

\usepackage{amsmath,amssymb,amsthm}
\usepackage{booktabs}
\usepackage{microtype}
\usepackage{url}

\newtheorem{theorem}{Theorem}
\newtheorem{lemma}[theorem]{Lemma}
\newtheorem{proposition}[theorem]{Proposition}
\newtheorem{corollary}[theorem]{Corollary}
\newtheorem{conjecture}[theorem]{Conjecture}
\theoremstyle{definition}
\newtheorem{definition}[theorem]{Definition}
\theoremstyle{remark}
\newtheorem{remark}[theorem]{Remark}

\title{The Character Structure\\of the Collision Fluctuation}

\author{Alexander S.\ Petty}

\email{alexander.petty@gmail.com}
\date{March 2023 (revised April 2026)}
\subjclass[2020]{11A63, 11N05, 11M41}

\begin{document}
\begin{abstract}
The truncated collision fluctuation sum
$\Phi_x(s, \ell) = \sum_{b < p \le x}
\Delta_p(\ell)\, p^{-s}$ appears to grow like
$-\mu_{\ell} \cdot \log\log x$ at $s = 1$.
We present four observations.
(1)~Computation in bases $3$, $6$, $7$, $10$, and $12$
is consistent with universality of this growth law.
(2)~The truncated sum decomposes over Dirichlet characters
modulo~$b$ (proved).
(3)~The character components do not appear to be
proportional to classical $L$-values (numerical evidence).
(4)~The per-frequency eigenvalue fluctuation sums appear
to converge to nearly the same value (numerical evidence).
\end{abstract}
\maketitle

% ============================================================
\section{Introduction}

This paper decomposes the collision fluctuation over Dirichlet characters.

The preceding paper~\cite{paper13} proved that the
collision fluctuation sum $\Phi_x(s, \ell)$ converges for
$\operatorname{Re}(s) > 2$ and conjectured logarithmic
divergence at $s = 1$:
\[
\Phi_x(1, \ell) = \sum_{p > b} \frac{\Delta_p(\ell)}{p}
\sim -\mu_{\ell} \cdot \log\log x.
\]
The present paper examines the structure of this divergence.

% ============================================================
\section{Universality Across Bases}

\begin{conjecture}[Universality]\label{conj:universal}
The Mertens growth law $\Phi_x(1, \ell) \sim -\mu_{\ell} \cdot
\log\log x$ and the convergence of $\Phi_x(s, \ell)$ for $s > 1$
are independent of the base~$b$.
\end{conjecture}

This conjecture is supported by computation in bases $3$, $6$,
$7$, $10$, and $12$, using all primes up to $8000$.

\begin{table}[h]
\centering
\begin{tabular}{rrrrr}
\toprule
base $b$ & $b - 1$ & $\Phi_x(1, 1)$ & $\Phi_x(0.9, 1)$ & behavior
  at $s = 0.9$ \\
\midrule
$3$ & $2$ & $-1.38$ & diverging & $\uparrow$ \\
$6$ & $5$ & $-1.29$ & diverging & $\uparrow$ \\
$7$ & $6$ & $-1.34$ & diverging & $\uparrow$ \\
$10$ & $9$ & $-1.37$ & diverging & $\uparrow$ \\
$12$ & $11$ & $-1.32$ & diverging & $\uparrow$ \\
\bottomrule
\end{tabular}
\caption{$\Phi_x(1,1)$ at $x = 8000$ grows logarithmically in every base tested.
$\Phi_x(0.9, 1)$ diverges as a power law. The critical boundary is
$s = 1$ in every base.}
\label{tab:bases}
\end{table}

The gate width ($b - 1$ derangements), the variance scaling
(standard deviation $\sim \sqrt{Q}$), and the critical boundary ($s = 1$)
are all base-independent. The base determines the quantitative
details (the Mertens constant~$\mu_{\ell}$ and the variance
constant~$\lambda_b$) but not the qualitative structure.

% ============================================================
\section{Decomposition over Dirichlet Characters}

The primes coprime to $b$ partition into $\phi(b)$ spectral
classes by their residue modulo~$b$. The Dirichlet characters
$\chi$ modulo~$b$ are the natural basis for functions on these
classes~\cite{davenport,iwaniec}.

\begin{definition}
The \emph{character-weighted fluctuation sum} for character
$\chi$ modulo~$b$ is
\[
L_{\chi}(s, \ell) = \sum_{p > b}
\frac{\chi(p)\, \Delta_p(\ell)}{p^s}.
\]
\end{definition}

\begin{proposition}[Character decomposition]
\label{prop:decomposition}
The fluctuation at primes in spectral class $r$ is
\[
F_r(s, \ell) = \sum_{\substack{p > b \\ p \equiv r \bmod b}}
\frac{\Delta_p(\ell)}{p^s}
= \frac{1}{\varphi(b)} \sum_{\chi}
\overline{\chi}(r)\, L_{\chi}(s, \ell).
\]
The total fluctuation sum is
$F = \sum_r F_r = L_{\chi_0}$, the trivial-character component.
\end{proposition}

\begin{proof}
Standard orthogonality of Dirichlet characters;
see~\cite{davenport}, Chapter~9, or~\cite{iwaniec}, Chapter~9.
\end{proof}

% ============================================================
\section{The Character Components}

\begin{conjecture}[Component structure]\label{conj:components}
Each character component $L_{\chi}(s, \ell)$ diverges
individually at $s = 1$ with a character-specific Mertens
constant.
\end{conjecture}

This is consistent with computation in base~$10$ at lag
$\ell = 1$ using all primes up to $10000$.

\begin{table}[h]
\centering
\begin{tabular}{rrrrrr}
\toprule
primes & $|L_{\chi_0}|$ & $|L_{\chi_1}|$ &
$|L_{\chi_2}|$ & $|L_{\chi_3}|$ \\
\midrule
$100$ & $1.073$ & $0.254$ & $0.321$ & $0.254$ \\
$500$ & $1.278$ & $0.356$ & $0.350$ & $0.356$ \\
$800$ & $1.341$ & $0.380$ & $0.355$ & $0.380$ \\
$1225$ & $1.393$ & $0.398$ & $0.350$ & $0.398$ \\
\bottomrule
\end{tabular}
\caption{Character component magnitudes at $s = 1$.
The trivial character $\chi_0$ carries the dominant weight.
$|L_{\chi_1}| = |L_{\chi_3}|$ by conjugate symmetry.}
\label{tab:components}
\end{table}

The trivial character $\chi_0$ carries the largest weight
($|L_{\chi_0}| \approx 1.4$). The non-trivial characters
contribute smaller but nonzero corrections
($|L_{\chi_j}| \approx 0.35$--$0.40$).

% ============================================================
\section{Relation to Classical $L$-Values}

The character components $L_{\chi}$ do not appear to be
proportional to classical Dirichlet $L$-values.

\begin{remark}
In base~$10$ at lag $\ell = 1$, the ratios
$L_{\chi_2}(1, 1) / L(1, \chi_2) \approx 0.54$ and
$|L_{\chi_1}(1, 1)| / |L(1, \chi_1)| \approx 0.40$
are not simple rational numbers and depend on the
lag~$\ell$. This suggests that the character components
$L_{\chi}(s, \ell)$ are not reducible to classical Dirichlet
$L$-functions~\cite{davenport,montgomery}.
\end{remark}

% ============================================================
\section{Spectral Equidistribution}

For primes $p$ where $b$ is a primitive root, the
cross-alignment matrix is permutation-similar to a
circulant with eigenvalues
$\lambda_k = \Phi(k) / L$, where $\Phi(k)$ is the spectral
power at frequency~$k$. As $p \to \infty$, the normalized
eigenvalue $\lambda_k / Q$ appears to converge to
\[
\lambda_k / Q \;\longrightarrow\;
\frac{b^2 \sin^2(\pi k / b)}{(\pi k)^2}.
\]
Define the eigenvalue fluctuation
$\psi_p(k) = \lambda_k / Q - b^2 \sin^2(\pi k/b) / (\pi k)^2$,
and the frequency fluctuation sum
\[
G(s, k) = \sum_{p > b} \frac{\psi_p(k)}{p^s}.
\]

\begin{conjecture}[Spectral equidistribution]\label{conj:spectral}
Each $G(1, k)$ converges at $s = 1$. Moreover, the limiting
values are approximately equal across all frequencies:
$G(1, k) \approx 0.165$ for $k = 1, \ldots, 5$ in base~$10$.
\end{conjecture}

\begin{table}[h]
\centering
\begin{tabular}{rrrrrr}
\toprule
primes & $G(1,1)$ & $G(1,2)$ & $G(1,3)$ & $G(1,4)$ & $G(1,5)$ \\
\midrule
$100$ & $0.1642$ & $0.1642$ & $0.1642$ & $0.1640$ & $0.1639$ \\
$500$ & $0.1650$ & $0.1650$ & $0.1649$ & $0.1646$ & $0.1644$ \\
$1000$ & $0.1651$ & $0.1651$ & $0.1650$ & $0.1646$ & $0.1644$ \\
\bottomrule
\end{tabular}
\caption{The per-frequency fluctuation sums converge to
nearly the same value for all $k$. The eigenvalue fluctuations
distribute uniformly across frequencies.}
\label{tab:spectral}
\end{table}

The near-equality of $G(1, k)$ across frequencies means that
the collision fluctuations are not concentrated at any
particular harmonic mode.

% ============================================================
\section{Polarity Equidistribution (Observational)}

The cumulative polarity of the repetend traces a path on the
fixed circle $\mathbb{Z}/(b{-}1)\mathbb{Z}$. At each
position of the repetend, the running digit sum modulo $b - 1$
takes a value in $\{1, \ldots, b{-}1\}$. The distribution of
these values is the \emph{polarity histogram}: the vector
$(h_1, \ldots, h_{b-1})$ where $h_j$ counts how many
repetend positions have running digit sum congruent to $j$
modulo $b - 1$.

For primes with $b$ as primitive root, the polarity histogram
appears to approach the uniform distribution $1/(b{-}1)$ as
$p \to \infty$. The path visits each value approximately
$L/(b{-}1)$ times. The polarity circle is the limiting
distribution: computation suggests the field approaches equidistribution
on it.

The collision fluctuations measure departure from this
limiting distribution. The Mertens growth law asserts that these departures,
summed over primes with weight $1/p$, accumulate at rate
$\log\log x$. The spectral equidistribution asserts that the
departures are flat across frequencies. Together: the field
approaches its equilibrium on the fixed polarity circle, and
the rate of approach is controlled by the critical boundary
$s = 1$.

% ============================================================
\section{Remarks}

\subsection*{The landscape}

The collision fluctuation now has four layers of structural
analysis: universality across bases, character decomposition
with independent components, new arithmetic objects distinct
from classical $L$-values, and spectral equidistribution of
eigenvalue fluctuations. The fluctuation decomposes along three
orthogonal axes (spectral class, frequency, polarity), and the
structure is consistent along each axis independently.

\subsection*{Open directions}

The character decomposition reduces the study of the Mertens
growth law to $\phi(b)$ independent problems: one for each
$L_{\chi}$ at $s = 1$. Whether any of these components admit
a closed-form Mertens constant, and whether the spectral
equidistribution has a theoretical explanation, are open
questions. The collision parameter
$c = b(1 - b^{\ell})^{-1} \bmod p$ (the residue class
that determines whether the multiplier $b^{\ell}$
deranges the bin partition) and its equidistribution
across spectral classes provide the analytic handle.

A modified fluctuation sum, corrected for the per-class
bias of~$S$, may converge at $s = 1$. If so, that centered
sum would be the collision analog of the prime number theorem.

% ============================================================
\begin{thebibliography}{9}

\bibitem{paper13}
A.~S.~Petty,
The collision fluctuation sum,
research note, 2024.

\bibitem{davenport}
H.~Davenport,
\emph{Multiplicative Number Theory},
3rd~ed., Springer, 2000.

\bibitem{iwaniec}
H.~Iwaniec and E.~Kowalski,
\emph{Analytic Number Theory},
AMS Colloquium Publications, vol.~53, 2004.

\bibitem{montgomery}
H.~L.~Montgomery and R.~C.~Vaughan,
\emph{Multiplicative Number Theory I: Classical Theory},
Cambridge University Press, 2007.

\end{thebibliography}

\end{document}
